\documentclass[11pt]{article}
\usepackage{amsmath,amssymb,amsthm,booktabs}
\usepackage[margin=1.2in]{geometry}
\usepackage{graphicx}
\usepackage{pgfplots}
\pgfplotsset{compat=1.18}
\usepackage[colorlinks=true,linkcolor=blue,citecolor=blue,urlcolor=blue]{hyperref}

\newtheorem{theorem}{Theorem}[section]
\newtheorem{proposition}[theorem]{Proposition}
\newtheorem{corollary}[theorem]{Corollary}
\newtheorem{lemma}[theorem]{Lemma}
\theoremstyle{remark}
\newtheorem{remark}[theorem]{Remark}
\newtheorem{example}[theorem]{Example}

\newcommand{\N}{\mathbb{N}}
\newcommand{\Prm}{\mathcal{P}}
\newcommand{\E}{\mathbb{E}}
\newcommand{\Prob}{\mathbb{P}}
\newcommand{\icv}{\preceq_{\mathrm{icv}}}

\title{Prime extremality for product integrals\\ over primitive sets}
\author{Alejandro Radisic\thanks{ORCID: \href{https://orcid.org/0009-0001-7612-1082}{0009-0001-7612-1082}.}}
\date{October 2026}

\emergencystretch=2em
\begin{document}
\maketitle

\begin{abstract}
A set $A$ of integers exceeding $1$ is primitive if no element divides
another. For $0<c\le1$ let $F_c(A)=\int_0^c\prod_{n\in A}(1-x^n)\,dx$.
We prove that among primitive subsets of $\{2,\dots,N\}$ the primes up
to $N$ are, for every $c$, the unique minimizer of $F_c$, and that over
all primitive sets the infimum is attained only by the set of all
primes, although the primes do not minimize the integrand pointwise. The
proof rests on one inequality,
$\int_0^c\prod_{t\in T}(1-x^t)\,dx>\int_0^c(1-x)\,dx$ for every set
$T\subseteq\{2,3,\dots\}$, which gives a strict compression replacing
the multiples of an integer $d$ by $d$. Consequently the
divisibility-minimal elements of any set $U$ uniquely minimize $F_c$
among primitive subsets of $U$. The minimum at $c=1$ is
$\Lambda=0.5506530112\ldots$. The truncation error satisfies
$-\log(I_N-\Lambda)\sim2\pi\sqrt{N/(3\log N)}$, with two further terms
from the prime number theorem, and to within $O(\log N)$ it is fixed by
the least prime exceeding $N$. The main results are formally verified
in Lean~4.
\end{abstract}

\noindent\textbf{Keywords.} primitive sets; divisibility antichains;
prime numbers; prime number theorem; partitions into primes; product
inequalities; stochastic orders; formal verification.

\section{Introduction}\label{sec:intro}

\subsection{Primitive sets and product integrals}
A set $A\subseteq\{2,3,\dots\}$ is \emph{primitive} (a divisibility
antichain) if no element of $A$ divides another. Write
$\Prm_N=\{p\le N:p\text{ prime}\}$ and $\Prm$ for the set of all
primes. To a set $A\subseteq\{2,3,\dots\}$, finite or infinite, we
attach the product
\[
  P_A(x)=\prod_{n\in A}(1-x^n)\qquad(0\le x<1),
\]
which converges locally uniformly on $[0,1)$ (for infinite $A$ because
$\sum_{n\in A}x^n<\infty$), and its \emph{initial integrals}
\[
  F_c(A)=\int_0^cP_A(x)\,dx,\qquad 0<c\le1 .
\]
For $A=\varnothing$ we set $P_\varnothing\equiv1$, so
$F_c(\varnothing)=c$.

Why this functional? The product $P_A$ is the reciprocal of the
generating function of partitions with parts in $A$,
\[
  \frac1{P_A(x)}=\sum_{n\ge0}p_A(n)\,x^n ,
\]
where $p_A(n)$ counts the partitions of $n$ into parts from $A$, and its
integral generalizes the classical evaluation of
$\int_0^1\prod_{n\ge1}(1-x^n)\,dx$ through Euler's pentagonal
theorem~\cite{Sandham1954,Mulcrone1956}. We ask which primitive set of
parts minimizes it. The answer is governed by a single comparison with
the one-part product $1-x$ (Theorem~\ref{thm:running} below).

Each $P_A$ is nonincreasing in $x$ and antitone in $A$: enlarging $A$
lowers $P_A$ pointwise. The extremal
question is therefore not whether large primitive sets beat small ones
but which \emph{maximal} antichain the integrals prefer. The answer is
the primes, uniformly in the truncation level $c$.

\begin{theorem}[Unique finite minimizer]\label{thm:unique}
Let $N\ge2$ and let $A\subseteq\{2,\dots,N\}$ be primitive with
$A\ne\Prm_N$. Then
\[
  F_c(\Prm_N)<F_c(A)\qquad\text{for every }c\in(0,1].
\]
\end{theorem}

\begin{theorem}[Infimum over all primitive sets; unique attainment]
\label{thm:profile}
Let $Q(x)=\prod_p(1-x^p)$. For every $c\in(0,1]$,
\[
  \inf_A F_c(A)=\int_0^cQ(x)\,dx=:\mathcal L(c),
\]
the infimum over \emph{all} primitive sets $A\subseteq\{2,3,\dots\}$,
finite or infinite, and it is attained precisely by $A=\Prm$: for every
primitive $A\ne\Prm$ and every $c\in(0,1]$, $F_c(A)>\mathcal L(c)$. In
particular no finite primitive set attains it, while the truncations
$\Prm_N$ form a minimizing sequence. The value at $c=1$ is
\[
  \Lambda:=\mathcal L(1)=\lim_{N\to\infty}\int_0^1\prod_{p\le N}(1-x^p)\,dx
  =0.5506530112728\ldots
\]
\end{theorem}

\subsection{Pointwise extremality fails}
The integrals are the right objects: the primes do \emph{not} minimize
the product itself. As $x\uparrow1$ one has
$P_A(x)\sim\bigl(\prod_{a\in A}a\bigr)(1-x)^{|A|}$ for finite $A$, so
near $1$ the product is governed by the cardinality of $A$, and maximal
antichains in $\{2,\dots,N\}$ can have about $N/2$ elements, far more
than $\pi(N)$. A concrete family already shows the failure at $N=9$.

\begin{proposition}[Failure of pointwise extremality]\label{prop:pointwise}
For $N\ge9$ let $A_N=(\Prm_N\setminus\{2,3\})\cup\{4,6,9\}$. Then $A_N$
is primitive and
\[
  \frac{P_{A_N}(x)}{P_{\Prm_N}(x)}=(1+x^2)(1+x^3)(1-x^9)
  \;\sim\;36\,(1-x)\longrightarrow0\qquad(x\uparrow1).
\]
Hence $P_{A_N}(x)<P_{\Prm_N}(x)$ on a neighbourhood of $1$, while
$F_c(\Prm_N)<F_c(A_N)$ for every $c\in(0,1]$ by
Theorem~\ref{thm:unique}.
\end{proposition}

\begin{proof}
No prime $\ge5$ divides or is divisible by $4$, $6$ or $9$, and these
three are pairwise incomparable, so $A_N$ is primitive. The ratio is
$(1-x^4)(1-x^6)(1-x^9)/\bigl((1-x^2)(1-x^3)\bigr)$.
\end{proof}

The ratio drops below $1$ at $x\approx0.965$. Thus every initial
integral prefers the primes although the integrand does not; the
theorems identify an integrated comparison that survives a genuine
obstruction to pointwise comparison (in probabilistic terms, the primes
are extremal in increasing-concave order but not in the usual stochastic
order; Remark~\ref{rem:icv}).

\subsection{The universal product inequality and prime compression}
Both theorems are iterations of a single operation. If $p$ is prime
and the primitive set $A$ contains multiples $pt$, $t\in T$, then
$p\notin A$, and the \emph{compression}
\[
  C_p(A)=\{p\}\cup\{a\in A:p\nmid a\}
\]
is again primitive. Writing $B=\{a\in A:p\nmid a\}$, the change in the
initial integral is
\begin{equation}\label{eq:exchange}
  F_c(A)-F_c(C_p(A))=\int_0^cP_B(u)\,E_T(u^p)\,du,
  \qquad
  E_T(v)=\prod_{t\in T}(1-v^t)-(1-v).
\end{equation}
For $|T|\ge2$ the factor $E_T$ changes sign on $[0,1]$ (already
$E_{\{2,3\}}(v)=v-v^2-v^3+v^5$), so the comparison cannot be pointwise. What is true is that its running
integral is positive, uniformly over all sets of exponents.

\begin{theorem}[Universal product inequality]\label{thm:running}
For every nonempty set $T\subseteq\{2,3,\dots\}$, finite or infinite,
and every $c\in(0,1]$,
\[
  \int_0^c\prod_{t\in T}(1-v^t)\,dv\;>\;\int_0^c(1-v)\,dv=c-\frac{c^2}{2}.
\]
\end{theorem}

Antitonicity reduces the theorem to $T=\{2,3,\dots\}$, where
$\prod_{n\ge2}(1-x^n)-(1-x)$ changes sign exactly once, so only the
endpoint $c=1$ must be checked (Section~\ref{sec:running}). The set
hypothesis is essential (Remark~\ref{rem:multiset}).
Section~\ref{sec:transport} transports the inequality through the
substitution $v=u^p$ and the weight $P_B$ in \eqref{eq:exchange},
giving the strict compression $F_c(C_p(A))<F_c(A)$, and
Section~\ref{sec:proofs} iterates it.

\subsection{Extremality relative to an ambient set}
Primality of $p$ plays no role in the compression inequality: the
substitution $v=u^d$ works for every integer $d\ge1$. What primality
secures is that $C_p(A)$ is again primitive. Replacing the primes by
arbitrary divisors gives a relative theorem. For
$U\subseteq\{2,3,\dots\}$ let
\[
  \operatorname{Min}(U)=\{u\in U:\ \text{no proper divisor of $u$ lies in $U$}\}
\]
be the set of divisibility-minimal elements of $U$. It is primitive, and
Theorem~\ref{thm:ambient} states that among primitive subsets of $U$ it is,
for every $c\in(0,1]$, the unique minimizer of $F_c$, for finite and
infinite sets alike. The case $U=\{2,3,\dots\}$ recovers
Theorem~\ref{thm:profile}. Two other cases correspond to recent work on
the additive weight $f(A)=\sum_{a\in A}1/(a\log a)$
(Corollary~\ref{cor:ambient}). For primitive sets contained in
$[x,\infty)$, the problem of Erd\H{o}s, S\'ark\"ozy and Szemer\'edi that
Alexeev et al.~\cite{Alexeev2026} settle asymptotically,
$f(A)\le1+O(1/\log x)$, has for $F_c$ an exact unique extremal set, the
integers $n\ge x$ whose largest proper divisor is less than $x$. For
primitive sets of integers with at least $k$ prime factors, all taken
from a set of primes $\mathcal Q$, the minimizer is the set of integers
with exactly $k$ such factors. This is the analogue of a conjecture of
Banks and Martin~\cite{BanksMartin2013}, and it holds for every
$\mathcal Q$, whereas the additive statement fails when
$2\in\mathcal Q$~\cite{Alexeev2026}.

\subsection{The truncations and the prime number theorem}
The minimizing sequence of Theorem~\ref{thm:profile} is
$I_N=F_1(\Prm_N)\downarrow\Lambda$. How fast it converges is a question
about the primes beyond $N$, and the answer is governed by their
distribution.

\begin{theorem}[Tail asymptotics]\label{thm:tail}
As $N\to\infty$,
\[
  -\log\bigl(I_N-\Lambda\bigr)\sim2\pi\sqrt{\frac{N}{3\log N}} .
\]
Consequently, if $p_n$ is the $n$th prime and $J_n=I_{p_n}$, then
$-\log(J_n-\Lambda)\sim(2\pi/\sqrt3)\sqrt n$.
\end{theorem}

With $\Phi(t)=-\log\prod_p(1-e^{-pt})$ and
$M_N=\inf_{t>0}(Nt+\Phi(t))$, the bound $I_N-\Lambda\le eN^2e^{-M_N}$ is
elementary and holds for every $N\ge3$. The prime number theorem is
needed exactly to evaluate $M_N$ and to place a prime in
$(N,(1+\varepsilon)N]$ (Section~\ref{sec:tailasymp}); Chebyshev's bounds give only the order of
magnitude (Remark~\ref{rem:chebyshev}). The leading term is also the
logarithmic growth rate of the number of partitions of $N$ into primes,
because both come from the same saddle point
(Remark~\ref{rem:partitions}).

Since $I_N$ changes only when $N$ passes a prime, the natural
variable is the least prime $q_N>N$. Elementary estimates show
$\bigl|{-\log(I_N-\Lambda)}-M_{q_N}\bigr|\le2\log q_N+4$
(Proposition~\ref{prop:nextprime}). Consequently, the bound
$-\log(I_N-\Lambda)\le M_N+O(\log N)$ is equivalent to the prime-gap
estimate $p_{n+1}-p_n\ll\sqrt{p_n}(\log p_n)^{3/2}$
(Corollary~\ref{cor:gaps}). That estimate follows from the Riemann
hypothesis but is not known unconditionally.

Relative lower-order terms, by contrast, are available unconditionally.
The prime number theorem with a classical error term gives
\[
  -\log(I_N-\Lambda)=2\pi\sqrt{\frac{N}{3\ell}}
  \Bigl(1-\frac{\log\ell}{2\ell}+\frac{C}{\ell}
  +O\Bigl(\frac{(\log\ell)^2}{\ell^2}\Bigr)\Bigr),\qquad \ell=\log N,
\]
with $C=\gamma-\zeta'(2)/\zeta(2)+\tfrac12\log(2\zeta(2))=1.7426\ldots$
(Theorem~\ref{thm:second}). Its remainder is much larger than $\log N$,
so this is consistent with Corollary~\ref{cor:gaps}.

\subsection{Context}
For the additive weight $\sum_{a\in A}1/(a\log a)$, extremality of the
primes over primitive sets is the Erd\H{o}s conjecture, proved by
Lichtman~\cite{Lichtman2023}; see Lichtman--Pomerance
\cite{LichtmanPomerance2019} for the history and \cite{Alexeev2026} for
recent developments. The present paper concerns a different, multiplicative
functional on the same class of sets, with an elementary proof; it is
not a new route to the additive problem. The comparison is
nonetheless instructive: for linear weights $\sum_{a\in A}w(a)$ the
primes win only with sufficient damping (for $w(a)=1/a$ they lose: the
integers in $[2,N]$ with exactly $k$ prime factors, $k\approx\log\log N$,
form a primitive set whose reciprocal sum is of order
$\log N/\sqrt{\log\log N}$; see \cite{LichtmanPomerance2019} and the
references there), whereas the product integrals prefer the
primes for every truncation level, with strict uniqueness and no
dependence on $N$, by a mechanism that is a single elementary
inequality.

Integrals of this shape have precedents. The all-exponent member
$\int_0^1\prod_{n\ge1}(1-x^n)\,dx$ is Sandham's Monthly Problem
4614~\cite{Sandham1954}, evaluated in closed form by
Mulcrone~\cite{Mulcrone1956} through Euler's pentagonal theorem
(\href{https://oeis.org/A258232}{OEIS A258232}); the even member is
\href{https://oeis.org/A258408}{OEIS A258408}. For the prime member no
closed form appears to be recorded. Its decimal expansion is
\href{https://oeis.org/A395518}{OEIS A395518}, contributed by the author, whose entry states the
finite-antichain infimum at $c=1$; the present paper supplies the
proofs, the uniformity in $c$, the strict uniqueness, and the infinite
case.

\section{The universal product inequality}\label{sec:running}

Throughout, $H(x)=\prod_{n\ge2}(1-x^n)$ for $0\le x<1$, and $H(1)=0$.
The product converges locally uniformly on $[0,1)$, and
$0\le H(x)\le1-x^2\to0$ as $x\uparrow1$, so $H$ is continuous on
$[0,1]$. For every set $T\subseteq\{2,3,\dots\}$ and $x\in[0,1]$,
\begin{equation}\label{eq:antitone}
  P_T(x)\ \ge\ H(x),
\end{equation}
since every factor omitted from $P_T$ lies in $[0,1]$. Two lemmas
control the function $H(x)-(1-x)$.

\begin{lemma}[Single crossing]\label{lem:crossing}
There is exactly one $x^*\in(0,1)$ with $H(x^*)=1-x^*$, and
$H(x)-(1-x)$ is strictly positive on $(0,x^*)$ and strictly negative on
$(x^*,1)$.
\end{lemma}

\begin{proof}
On $[0,1)$ write
\[
  h(x)=\log\frac{H(x)}{1-x}=\log(1+x)+\sum_{n\ge3}\log(1-x^n),
\]
using $(1-x^2)/(1-x)=1+x$. The series converges, together with the
series of first and second derivatives, locally uniformly on $[0,1)$
(each is dominated by a multiple of $\sum n^2x^{n-2}$ on $[0,r]$,
$r<1$), so $h$ is twice differentiable there with
\[
  h''(x)=-\frac1{(1+x)^2}-\sum_{n\ge3}
  \frac{n\,x^{n-2}\bigl((n-1)+x^n\bigr)}{(1-x^n)^2}\ <\ 0 .
\]
Thus $h$ is strictly concave on $[0,1)$. Moreover $h(0)=0$,
$h'(0)=1>0$, and $h(x)\le\log(1+x)+\log(1-x^3)\to-\infty$ as
$x\uparrow1$. A strictly concave function that vanishes at $0$ and
increases there is positive on some interval $(0,\delta)$, and the set
where it is nonnegative is an interval; since $h\to-\infty$, that
interval is $[0,x^*]$ for a unique $x^*\in(0,1)$, with $h>0$ on
$(0,x^*)$ and $h<0$ on $(x^*,1)$. Since
$H(x)-(1-x)=(1-x)(e^{h(x)}-1)$ has the sign of $h$ on $(0,1)$, the lemma
follows.
\end{proof}

\begin{lemma}[Endpoint estimate]\label{lem:endpoint}
\[
  \int_0^1H(x)\,dx-\frac12\ \ge\ \frac{34619}{4476780}\ >\ 0 .
\]
\end{lemma}

\begin{proof}
For numbers $y_n\in[0,1]$ one has $\prod(1-y_n)\ge1-\sum y_n$ (the
inequality is trivial when the right side is negative and follows by
induction otherwise). Applying this to the factors with $n\ge6$ and
keeping the factors $2,\dots,5$ exactly gives, for $0\le x<1$,
\[
  H(x)\ \ge\ \prod_{n=2}^{5}(1-x^n)\Bigl(1-\frac{x^6}{1-x}\Bigr)
  =(1+x)(1-x^3)(1-x^4)(1-x^5)\,(1-x-x^6),
\]
a polynomial of degree $19$ (the bound may be negative near $1$, which
does no harm). Exact integration gives
\[
  \int_0^1\prod_{n=2}^{5}(1-x^n)\,dx=\frac{96601}{180180},
  \qquad
  \int_0^1\frac{x^6\prod_{n=2}^{5}(1-x^n)}{1-x}\,dx=\frac{275501}{9699690},
\]
using $\prod_{n=2}^{5}(1-x^n)
=1-x^2-x^3-x^4+x^6+2x^7+x^8-x^{10}-x^{11}-x^{12}+x^{14}$ and
$(1+x)(1-x^3)(1-x^4)(1-x^5)
=1+x-x^3-2x^4-2x^5-x^6+x^7+2x^8+2x^9+x^{10}-x^{12}-x^{13}$.
The difference of the two integrals minus $\tfrac12$ is
$34619/4476780$.
\end{proof}

The true value is $\int_0^1H\,dx-\tfrac12=0.01609\ldots$, so the
estimate loses only a factor of about two; the crossing is at
$x^*=0.567\ldots$ (Figure~\ref{fig:crossing}).

\begin{proof}[Proof of Theorem~\ref{thm:running}]
Let $R(c)=\int_0^c\bigl(H(x)-(1-x)\bigr)dx$. By
Lemma~\ref{lem:crossing}, $R$ is strictly increasing on $[0,x^*]$ and
strictly decreasing on $[x^*,1]$, with $R(0)=0$. Hence $R(c)>0$ for
$c\in(0,x^*]$, and for $c\in[x^*,1]$, $R(c)\ge R(1)>0$ by
Lemma~\ref{lem:endpoint}. So $R>0$ on $(0,1]$. By
\eqref{eq:antitone}, $\int_0^cP_T\ge\int_0^cH$ for every $T$, and the
theorem follows.
\end{proof}

\begin{remark}[Sets versus multisets]\label{rem:multiset}
The function $H$ is the pointwise infimum of $P_T$ over \emph{sets}
$T$, and that is exactly what \eqref{eq:antitone} uses. For multisets
the inequality fails: with three copies of the exponent $2$,
$\int_0^1(1-v^2)^3\,dv=\tfrac{16}{35}<\tfrac12$, and indeed
$\int_0^c(1-v^2)^3\,dv<c-c^2/2$ for all $c>0.6451\ldots$. Multisets
never arise below because $A$ is a set.
\end{remark}

\begin{figure}[t]
\centering
\begin{tikzpicture}
\begin{axis}[
  width=0.62\textwidth, height=0.40\textwidth,
  xlabel={$x$}, xmin=0, xmax=1,
  ymin=-0.21, ymax=0.21, ytick={-0.2,-0.1,0,0.1,0.2},
  yticklabel style={/pgf/number format/fixed},
  legend style={draw=none, font=\small, at={(0.5,1.03)}, anchor=south,
    legend columns=2, /tikz/every even column/.append style={column sep=1.5em}},
  axis lines=left, tick label style={font=\small},
  label style={font=\small},
]
\addplot[black!50, thin, forget plot] coordinates {(0,0) (1,0)};
\addplot[blue, thick] table[x=x, y=E] {figures/crossing.dat};
\addlegendentry{$H(x)-(1-x)$}
\addplot[red, thick, dashed] table[x=x, y=G] {figures/crossing.dat};
\addlegendentry{$\int_0^x\bigl(H-(1-u)\bigr)\,du$}
\end{axis}
\end{tikzpicture}
\caption{The extremal case $T=\{2,3,\dots\}$ of
Theorem~\ref{thm:running}. The integrand $H(x)-(1-x)$ crosses zero
exactly once, at $x^*\approx0.567$, so its running integral rises and
then falls to the endpoint value $\int_0^1H-\tfrac12\approx0.0161$,
which is positive.}
\label{fig:crossing}
\end{figure}

\section{Transport and strict compression}\label{sec:transport}

The block exchange \eqref{eq:exchange} needs
Theorem~\ref{thm:running} after the substitution $v=u^p$ and under a
nonincreasing weight $P_B$. The lemmas below transport the estimate,
and Proposition~\ref{prop:block} combines them.

\begin{lemma}[Singular monotone weights]\label{lem:weight}
Let $g$ be continuous on $[0,1]$ with $G(y)=\int_0^yg>0$ for all
$y\in(0,1]$, and let $w$ be positive, nonincreasing and $C^1$ on
$(0,1]$, integrable at $0$, with $w(\varepsilon)G(\varepsilon)\to0$
as $\varepsilon\downarrow0$. Then for every $c\in(0,1]$,
\[
  \int_0^cw(v)g(v)\,dv\ \ge\ w(c)\,G(c)\ >\ 0 .
\]
\end{lemma}

\begin{proof}
Integrate by parts on $[\varepsilon,c]$:
$\int_\varepsilon^cwg=w(c)G(c)-w(\varepsilon)G(\varepsilon)
-\int_\varepsilon^cw'G$. The last integral is nonpositive since
$w'\le0$ and $G\ge0$, and $w(\varepsilon)G(\varepsilon)\to0$; let
$\varepsilon\downarrow0$ (the left side converges since $wg$ is
integrable at $0$).
\end{proof}

\begin{corollary}[Power substitution]\label{lem:power}
Let $p\ge1$ be an integer and let $T\subseteq\{2,3,\dots\}$ be
nonempty, finite or infinite. Then for every $c\in(0,1]$,
\[
  \int_0^cE_T(u^p)\,du>0 .
\]
\end{corollary}

\begin{proof}
Let $G(y)=\int_0^yE_T(v)\,dv$, so $G>0$ on $(0,1]$ by
Theorem~\ref{thm:running}, and $G(y)=O(y)$ at $0$ since $E_T$ is
continuous on $[0,1]$. Substituting $v=u^p$,
\[
  \int_0^cE_T(u^p)\,du=\frac1p\int_0^{c^p}v^{1/p-1}E_T(v)\,dv,
\]
and $w(v)=v^{1/p-1}/p$ is positive, nonincreasing, integrable at $0$,
with $w(\varepsilon)G(\varepsilon)=O(\varepsilon^{1/p})\to0$; apply
Lemma~\ref{lem:weight}.
\end{proof}

\begin{lemma}[Monotone weights]\label{lem:stieltjes}
Let $g$ be continuous on $[0,1]$ with $G(y)=\int_0^yg\ge0$ for all
$y\in[0,1]$, and let $w$ be nonnegative, continuous and nonincreasing on
$[0,1]$. Then for every $c\in(0,1]$,
\[
  \int_0^cw(u)g(u)\,du
  =w(c)G(c)+\int_{(0,c]}G\,d(-w)\ \ge\ 0,
\]
where $d(-w)$ is the finite nonnegative Borel measure on $(0,c]$ with
$d(-w)\bigl((a,b]\bigr)=w(a)-w(b)$. The left side is strictly positive
if $G>0$ on $(0,c]$ and either $w(c)>0$ or $d(-w)$ has positive mass on
$(0,c)$.
\end{lemma}

\begin{proof}
Since $w$ is continuous and nonincreasing, $d(-w)$ is a finite
nonnegative measure without atoms, of total mass $w(0)-w(c)$. As
$G'=g$ is continuous, $\int_0^cwg\,du=\int_0^cw\,dG$, and
Riemann--Stieltjes integration by parts with $G(0)=0$ gives the stated
identity; both terms are nonnegative. For strictness: if $w(c)>0$ then
$w(c)G(c)>0$; if $d(-w)$ has positive mass on $(0,c)$, then, $G$ being
continuous and strictly positive there, $\int G\,d(-w)>0$.
\end{proof}

The Stieltjes form is needed because the weight will be an infinite
product $P_B$, which is continuous and nonincreasing on $[0,1]$ but need
not be $C^1$ at the endpoint; Lemma~\ref{lem:weight}, by contrast,
allows a singularity at $0$ but needs a smooth weight.

\begin{proposition}[Strict block exchange]\label{prop:block}
Let $p\ge1$ be an integer (primality is not needed) and let
$B,T\subseteq\{2,3,\dots\}$ be finite or
infinite, $T$ nonempty, with $B$ disjoint from $pT=\{pt:t\in T\}$ and
$p\notin B$. Set $A=B\cup pT$ and $A'=B\cup\{p\}$. Then for every
$c\in(0,1]$,
\[
  F_c(A')<F_c(A).
\]
\end{proposition}

\begin{proof}
Pointwise,
$P_A(u)-P_{A'}(u)=P_B(u)\bigl(\prod_{t\in T}(1-u^{pt})-(1-u^p)\bigr)
=P_B(u)E_T(u^p)$, so
\[
  F_c(A)-F_c(A')=\int_0^cP_B(u)\,E_T(u^p)\,du .
\]
The weight $w=P_B$ is nonnegative, continuous and nonincreasing on
$[0,1]$, and $G(y)=\int_0^yE_T(u^p)\,du>0$ on $(0,1]$ by
Corollary~\ref{lem:power}. If $c<1$ then $P_B(c)>0$ and
Lemma~\ref{lem:stieltjes} gives strict positivity. If $c=1$: either
$B=\varnothing$, and the integral is $\int_0^1E_T(u^p)\,du>0$ directly;
or $B\ne\varnothing$, in which case $P_B(1)=0$ and $d(-P_B)$ has total
mass $P_B(0)-P_B(1)=1$, carried by $(0,1)$ since $P_B$ is continuous,
while $G>0$ there; the strictness clause of Lemma~\ref{lem:stieltjes}
applies.
\end{proof}

\section{Proofs of the extremality theorems}\label{sec:proofs}

For a primitive $A$ containing a composite $m$ and a prime $p\mid m$,
primitivity gives $p\notin A$, and the compression
$C_p(A)=\{p\}\cup\{a\in A:p\nmid a\}$ is a block exchange in the sense
of Proposition~\ref{prop:block} with $B=\{a\in A:p\nmid a\}$ and
nonempty $T=\{a/p:a\in A,\ p\mid a\}$. It preserves primitivity:
inherited nonmultiples remain incomparable, $p$ divides none of them,
and none divides $p$.

\begin{proof}[Proof of Theorem~\ref{thm:unique}]
Let $A\subseteq\{2,\dots,N\}$ be primitive, $A\ne\Prm_N$, and fix
$c\in(0,1]$. If $A$ contains a composite $m$, choose a prime $p\mid m$;
then $F_c(C_p(A))<F_c(A)$ by Proposition~\ref{prop:block}. The
compression strictly lowers the element sum, since a nonempty block of
multiples of $p$, all at least $2p$, is replaced by $p$. Induction on
the element sum terminates at an all-prime set $R\subseteq\Prm_N$ with
$F_c(R)<F_c(A)$, strict inequality being secured by the first
compression step.

If instead $A$ contains no composite, then $A\subseteq\Prm_N$ with
$A\ne\Prm_N$, and we take $R=A$. In either case, if $R\ne\Prm_N$, add
the missing primes one at a time: for $q\in\Prm_N\setminus R$,
\[
  F_c(R\cup\{q\})=F_c(R)-\int_0^cx^qP_R(x)\,dx<F_c(R),
\]
since the integrand is strictly positive on $(0,c)$. Chaining,
$F_c(\Prm_N)<F_c(A)$.
\end{proof}

\begin{proof}[Proof of Theorem~\ref{thm:profile}]
Fix $c\in(0,1]$.

\emph{Lower bound.} For finite primitive $A\subseteq\{2,\dots,N\}$,
Theorem~\ref{thm:unique} (with its trivial case $A=\Prm_N$) gives
$F_c(A)\ge F_c(\Prm_N)$. The integrands $P_{\Prm_N}$ decrease pointwise
to $Q$ and are bounded by $1$, so by dominated convergence
$F_c(\Prm_N)\downarrow\int_0^cQ=\mathcal L(c)$; hence
$F_c(A)\ge\mathcal L(c)$ for every finite primitive $A$. For infinite
primitive $A$, the truncations $A\cap[2,N]$ are finite primitive and
$F_c(A)=\lim_NF_c(A\cap[2,N])\ge\mathcal L(c)$, again by dominated
convergence. Attainment at the primes is the definition:
$F_c(\Prm)=\int_0^cQ=\mathcal L(c)$.

\emph{Uniqueness.} Let $A\ne\Prm$ be primitive; we show
$F_c(A)>\mathcal L(c)$. If $A=\varnothing$ then $F_c(A)=c>\mathcal L(c)$
since $Q<1$ on $(0,c)$. If $A$ contains a composite $m$, choose a prime
$p\mid m$; the compression $C_p(A)$ is a block exchange with nonempty
$T$ (possibly infinite $T$ and $B$), preserves primitivity, and by
Proposition~\ref{prop:block} satisfies $F_c(C_p(A))<F_c(A)$.
Therefore $F_c(A)>F_c(C_p(A))\ge\mathcal L(c)$, the second inequality
being the lower bound applied to the primitive set $C_p(A)$. If $A$
contains no composite, then $A\subsetneq\Prm$; choose a prime
$q\notin A$. Then $A\cup\{q\}$ is primitive and
\[
  F_c(A)-F_c(A\cup\{q\})=\int_0^cx^qP_A(x)\,dx>0,
\]
since $P_A>0$ on $(0,1)$, so
$F_c(A)>F_c(A\cup\{q\})\ge\mathcal L(c)$.

\emph{Non-attainment by finite sets.} Every finite primitive set differs
from $\Prm$, so the strict inequality above proves non-attainment; the
minimizing sequence $\Prm_N$ was exhibited in the lower bound. The
decimal expansion of $\Lambda=\mathcal L(1)$ is certified in
Section~\ref{sec:Lambda}.
\end{proof}

\subsection{Extremality relative to an ambient set}

For $U\subseteq\{2,3,\dots\}$ let $\operatorname{Min}(U)$ be the set of
$u\in U$ with no proper divisor in $U$. It is primitive: if $a\mid b$
with $a\ne b$ both in $\operatorname{Min}(U)$, then $a$ is a proper
divisor of $b$ lying in $U$.

\begin{theorem}[Relative extremality]\label{thm:ambient}
Let $U\subseteq\{2,3,\dots\}$ and let $A\subseteq U$ be primitive, finite
or infinite, with $A\ne\operatorname{Min}(U)$. Then
\[
  F_c(\operatorname{Min}(U))<F_c(A)\qquad\text{for every }c\in(0,1].
\]
Likewise, for every $N$ and every primitive $A\subseteq U\cap[2,N]$ with
$A\ne\operatorname{Min}(U)\cap[2,N]$,
$F_c(\operatorname{Min}(U)\cap[2,N])<F_c(A)$ for every $c\in(0,1]$.
\end{theorem}

\begin{proof}
Let $A\subseteq U$ be primitive and let $a\in A\setminus\operatorname{Min}(U)$.
Choose a proper divisor $d\in U$ of $a$. No element $b\in A$ divides $d$:
otherwise $b\mid a$ and so $b=a$ by primitivity, whence $a\mid d\mid a$
and $d=a$. In particular $d\notin A$. Hence
$C_d(A)=\{d\}\cup\{b\in A:d\nmid b\}$ is primitive, as in
Section~\ref{sec:proofs}, and $C_d(A)\subseteq U$. Moreover
$F_c(C_d(A))<F_c(A)$ by Proposition~\ref{prop:block}, applied with $p=d$,
$B=\{b\in A:d\nmid b\}$ and $T=\{b/d:b\in A,\ d\mid b\}\ni a/d$.
Every element of $T$ is at least $2$ because $d\notin A$.

The proofs of Theorems~\ref{thm:unique} and~\ref{thm:profile} now go
through verbatim, with the compressions $C_d$ in place of $C_p$ and
$\operatorname{Min}(U)$ in place of the primes. In the finite case,
compression strictly lowers the element sum, and a set contained in
$\operatorname{Min}(U)\cap[2,N]$ is completed by adding the missing
elements one at a time, each addition strictly lowering $F_c$. For the
lower bound in the infinite case, the truncations $A\cap[2,N]$ are
primitive subsets of $U\cap[2,N]$, and
$\operatorname{Min}(U\cap[2,N])=\operatorname{Min}(U)\cap[2,N]$ because
divisors are smaller than their multiples; dominated convergence then
gives $F_c(\operatorname{Min}(U))\le F_c(A)$. For strictness, either
$A\subsetneq\operatorname{Min}(U)$, and adding a missing element of
$\operatorname{Min}(U)$ keeps the set primitive and strictly lowers
$F_c$; or $A$ has an element outside $\operatorname{Min}(U)$, and one
compression $C_d$ strictly lowers $F_c$. In both cases the lower bound,
applied to the new set, finishes the proof.
\end{proof}

For $U=\{2,3,\dots\}$, $\operatorname{Min}(U)$ is the set of primes, and
Theorem~\ref{thm:ambient} contains Theorems~\ref{thm:unique}
and~\ref{thm:profile}. Write $\Omega(n)$ for the number of prime factors
of $n$ counted with multiplicity and $P^-(n)$ for the least prime
factor.

\begin{corollary}\label{cor:ambient}
For every $c\in(0,1]$ the following sets are the unique minimizers of
$F_c$ in the stated classes.
\begin{enumerate}
\item[(a)] \emph{Large numbers.} Among primitive sets contained in
$[x,\infty)$, $x\ge2$: the set
$S_x=\{n\ge x:\ n/P^-(n)<x\}$ of integers $n\ge x$ whose largest proper
divisor is less than $x$. It contains $[x,2x)$ and every prime $\ge x$,
and $S_2=\Prm$.
\item[(b)] \emph{Restricted primes.} Let $\mathcal Q$ be a nonempty set
of primes, let $\N_k(\mathcal Q)$ be the set of integers with
$\Omega(n)=k$ all of whose prime factors lie in $\mathcal Q$, and
$\N_{\ge k}(\mathcal Q)=\bigcup_{j\ge k}\N_j(\mathcal Q)$. Among
primitive subsets of $\N_{\ge k}(\mathcal Q)$: the set
$\N_k(\mathcal Q)$. In particular
\[
  F_c(\N_1(\mathcal Q))<F_c(\N_2(\mathcal Q))<F_c(\N_3(\mathcal Q))<\cdots .
\]
\item[(c)] \emph{Hilbert numbers.} Among primitive sets of integers
$n>1$ with $n\equiv1\pmod4$: the irreducible elements of the
multiplicative monoid $1+4\N$ (for instance $5,9,13,21,49$).
\end{enumerate}
\end{corollary}

\begin{proof}
In each case we identify $\operatorname{Min}(U)$ and apply
Theorem~\ref{thm:ambient}. (a) A proper divisor of $n$ lies in
$[x,\infty)$ if and only if the largest one, $n/P^-(n)$, is at least
$x$. For $n<2x$ it is at most $n/2<x$. (b) If $\Omega(n)>k$, deleting
prime factors of $n$ produces a proper divisor in $\N_k(\mathcal Q)$; if
$\Omega(n)=k$, every proper divisor has fewer than $k$ prime factors. The
chain follows because $\N_{k+1}(\mathcal Q)$ is a primitive subset of
$\N_{\ge k}(\mathcal Q)$ different from the nonempty set
$\N_k(\mathcal Q)$. (c) This is the definition of irreducibility in the
monoid.
\end{proof}

\begin{remark}[Comparison with the additive weight]\label{rem:additive}
For $f(A)=\sum_{a\in A}1/(a\log a)$, the problem of
Corollary~\ref{cor:ambient}(a) is a 1966 question of Erd\H{o}s,
S\'ark\"ozy and Szemer\'edi. Alexeev et al.~\cite{Alexeev2026} prove
$f(A)\le1+O(1/\log x)$ for primitive $A\subseteq[x,\infty)$, by a Markov
chain with von Mangoldt weights that moves from $n$ to $n/q$. The
analogue of (b) for $f$ is a conjecture of Banks and
Martin~\cite{BanksMartin2013}; it is false when $2\in\mathcal Q$ and is
proved in~\cite{Alexeev2026} for sets of odd primes. For the product
integrals both statements are exact and hold without restriction. The
reason is that the compression inequality behind
Theorem~\ref{thm:ambient} holds uniformly in the divisor $d$ and in the
block of multiples $dT$, while linear weights must balance the lost
multiples against the single gained divisor.
\end{remark}

\begin{remark}[Stochastic order and moments]\label{rem:icv}
Let $M_n$ be the maximum of $n$ independent uniform $[0,1]$ variables,
the families independent as $n$ varies, and let $Z_A=\inf_{n\in A}M_n$
($Z_\varnothing=1$). By Borel--Cantelli the infimum is attained almost
surely for infinite $A$, so $\Prob(Z_A>x)=P_A(x)$ and
$F_c(A)=\E\min(Z_A,c)$. The hinge functions $\min(\,\cdot\,,c)$
generate the increasing-concave order $\icv$~\cite{ShakedShanthikumar}:
after replacing $\phi(0)$ by $\phi(0^+)$, which changes no expectation
since $Z_A>0$ almost surely, every increasing concave $\phi$ on $[0,1]$
has the form $\phi(0^+)+\alpha x+\int_{(0,1]}\min(x,t)\,\mu(dt)$ with
$\alpha\ge0$ and $\mu\ge0$. Theorems~\ref{thm:unique}
and~\ref{thm:profile} therefore give $Z_{\Prm_N}\icv Z_A$ for primitive
$A\subseteq\{2,\dots,N\}$ and $Z_\Prm\icv Z_A$ for every primitive $A$,
whereas by Proposition~\ref{prop:pointwise} the usual stochastic order
fails.

The extremality also reaches beyond concave test functions: for
$0<k\le17/4$ the primes uniquely minimize
$\E\min(Z_A^k,c^k)=\int_0^cku^{k-1}P_A(u)\,du$, for finite and infinite
primitive sets alike; for instance $x+\varepsilon x^2$ singles out the
primes for every $\varepsilon>0$. A complete proof is formalized
(Section~\ref{sec:formal}); we sketch it. For $k\le1$ it follows from
the case $k=1$ by integration by parts. For $k>1$ one repeats
Sections~\ref{sec:transport} and~\ref{sec:proofs} with the weight
$ku^{k-1}$. A compression at $p\ge k$ is handled by
Lemma~\ref{lem:weight}. For $p<k$, put $s=k/p$ and
$J_s(z)=\int_0^zx^{s-1}\bigl(Q(x)-(1-x)\bigr)dx$; the change in the
integral splits into the corresponding comparison of the quotient block
$T=\{a/p\}$ with the primes at exponent $s$, which is nonnegative by
induction on $k$ because $T$ is itself primitive, plus $J_s(y^p)$. As in
Lemma~\ref{lem:crossing}, $Q(x)-(1-x)$ changes sign once, so
$J_s(z)>0$ for all $z$ once $J_s(1)>0$, and $J_{17/8}(1)>1/20000$ is an
exact rational computation from
$Q(x)\ge(1-x^2-x^{19})\prod_{3\le p\le17}(1-x^p)$ (script
\texttt{moment\_certificate.py}). Three rounds give
$[1,2]\to[1,4]\to[1,17/4]$. The method stops at
$2\alpha_0\approx4.286$, where $\alpha_0=\sup\{b:J_b(1)>0\}$
(numerically; $J_3(1)<0$ exactly). Some restriction on $k$ is
necessary (open problem~2).
\end{remark}

\section{The limiting value and its truncations}\label{sec:consequences}

\subsection{The limiting value}\label{sec:Lambda}

For $c_0<1$ the series $\sum_pc_0^p$ converges, so the product defining
$Q$ converges uniformly on $[0,c_0]$; hence $Q$ is continuous on $[0,1)$
and $\mathcal L$ is differentiable on $(0,1)$ with
$\mathcal L'(c)=Q(c)$. Write $I_N=F_1(\Prm_N)$, so that
$I_N\downarrow\Lambda=\mathcal L(1)$. An exact rational sandwich gives
\[
  0.5506530112728429635779717568\ldots<\Lambda
  <0.5506530112728429635779717571\ldots,
\]
certifying $26$ decimal places. The upper end is $I_{2003}$, computed
exactly by integrating the integer coefficient vector of
$\prod_{p\le2003}(1-x^p)$ (degree $279{,}053$) termwise. For the lower
end, every prime $p>2003$ is odd and at least $2005$, so
$1-\prod_{p>2003}(1-x^p)\le x^{2005}/(1-x^2)$ by the product inequality
of Lemma~\ref{lem:endpoint}, and we subtract
$\int_0^1x^{2005}\prod_{2<p\le2003}(1-x^p)\,dx$.

\subsection{Tail asymptotics}\label{sec:tailasymp}
Throughout, $p$ runs over primes, $a=\zeta(2)=\pi^2/6$,
$L(t)=\log(1/t)$, and for $u,t>0$
\[
  S(u)=\sum_pe^{-pu},\qquad W(u)=\sum_p\log p\,e^{-pu},\qquad
  \Phi(t)=\sum_p-\log\bigl(1-e^{-pt}\bigr),
\]
so that $Q(e^{-t})=e^{-\Phi(t)}$. Each summand of $\Phi$ is positive and
decreasing in $t$, so $\Phi$ is decreasing. We use the prime number
theorem in the form $\theta(x)=\sum_{p\le x}\log p=x+o(x)$.

\begin{lemma}[Exponential prime sums]\label{lem:expsums}
As $u\downarrow0$, $uW(u)\to1$ and $uL(u)S(u)\to1$.
\end{lemma}

\begin{proof}
Partial summation gives
\[
  W(u)=\sum_{n\ge1}(\theta(n)-\theta(n-1))e^{-nu}
  =(1-e^{-u})\sum_{n\ge1}\theta(n)e^{-nu},
\]
the series converging since $\theta(n)\le n\log4$. Given $\varepsilon>0$ there is $C_\varepsilon$
with $|\theta(n)-n|\le\varepsilon n+C_\varepsilon$ for all $n$. Since
$\sum_nne^{-nu}=e^{-u}/(1-e^{-u})^2$, we get
\[
  \Bigl|W(u)-\frac{e^{-u}}{1-e^{-u}}\Bigr|
  \le\frac{\varepsilon e^{-u}}{1-e^{-u}}+C_\varepsilon .
\]
Multiplying by $u$ and letting $u\downarrow0$ gives $uW(u)\to1$.

For $S$, fix $\delta\in(0,1)$. Primes $p\le u^{-(1-\delta)}$ contribute at
most $u^{-(1-\delta)}=o(1/(uL(u)))$ to $S(u)$. The remaining primes have
$\log p\ge(1-\delta)L(u)$, so they contribute at most
$W(u)/((1-\delta)L(u))$. Hence $\limsup uL(u)S(u)\le1/(1-\delta)$.
Conversely, primes $p\le Y=u^{-(1+\delta)}$ have $\log p\le(1+\delta)L(u)$,
so
\[
  (1+\delta)L(u)S(u)\ \ge\ W(u)-\sum_{n>Y}\log n\,e^{-nu}
  \ \ge\ W(u)-e^{-Yu/2}\sum_{n\ge1}\log n\,e^{-nu/2}.
\]
Here $Yu=u^{-\delta}\to\infty$ and the last sum is $O(u^{-2})$, so the
subtracted term is $o(1/u)$ and $\liminf uL(u)S(u)\ge1/(1+\delta)$. Let
$\delta\downarrow0$.
\end{proof}

\begin{lemma}[The prime partition exponent]\label{lem:phi}
As $t\downarrow0$, $tL(t)\Phi(t)\to\zeta(2)$.
\end{lemma}

\begin{proof}
Expanding $-\log(1-y)=\sum_{m\ge1}y^m/m$ and exchanging the order of
summation (all terms are nonnegative) gives
$\Phi(t)=\sum_{m\ge1}S(mt)/m$. Since $L(t)=L(mt)+\log m$, the $m$th term
of $tL(t)\Phi(t)$ is
\[
  \frac{tL(t)S(mt)}{m}
  =\frac{1}{m^2}\,mt\,L(mt)S(mt)\Bigl(1+\frac{\log m}{L(mt)}\Bigr),
\]
which tends to $1/m^2$ for each fixed $m$ by Lemma~\ref{lem:expsums}.
For domination, $S(u)\le\sum_{n\ge1}e^{-nu}\le1/u$ for every $u>0$, and
by Lemma~\ref{lem:expsums} there is $u_1\in(0,e^{-1}]$ with
$uL(u)S(u)\le2$ for $0<u<u_1$. If $mt<u_1$, then $L(mt)\ge1$ and the $m$th
term is at most $2(1+\log m)/m^2$. If $mt\ge u_1$, it is at most
$L(t)/m^2=(L(mt)+\log m)/m^2\le(L(u_1)+\log m)/m^2$. The bound
$(2+L(u_1))(1+\log m)/m^2$ is summable, so dominated convergence gives
$tL(t)\Phi(t)\to\sum_m1/m^2=\zeta(2)$.
\end{proof}

Lemma~\ref{lem:phi} is the classical asymptotic for the logarithm of
the generating function of partitions into primes; we include the short
proof because the formalization follows it. The next lemma is
unconditional. It compares $I_N-\Lambda$ with the function $\Phi$,
setting
\[
  M_N=\inf_{t>0}\bigl(Nt+\Phi(t)\bigr).
\]

\begin{lemma}[Sandwich]\label{lem:sandwich}
\begin{enumerate}
\item[(a)] For every $N\ge3$, $0\le I_N-\Lambda\le e\,N^2e^{-M_N}$.
\item[(b)] Let $q>N$ be prime, $\tau>0$ and $\varepsilon>0$. If
$\Phi(t)\le B$ for $\tau\le t\le(1+\varepsilon)\tau$, then
\[
  I_N-\Lambda\ \ge\ \frac{\varepsilon\tau e^{-\tau}}{1+\varepsilon\tau}
  \exp\bigl(-B-(1+\varepsilon)q\tau\bigr).
\]
\end{enumerate}
\end{lemma}

\begin{proof}
Write $P_N=P_{\Prm_N}$. For $M\ge N$ and $0\le x<1$,
\[
  0\le P_N(x)-P_M(x)=P_N(x)\Bigl(1-\prod_{N<p\le M}(1-x^p)\Bigr)
  \le P_N(x)\min\Bigl(1,\frac{x^{N+1}}{1-x}\Bigr),
\]
by the product inequality of Lemma~\ref{lem:endpoint} and
$\sum_{n>N}x^n=x^{N+1}/(1-x)$. Integrating and letting $M\to\infty$
($I_M\to\Lambda$) gives $I_N-\Lambda\le\sup_{0<x<1}g(x)$, where $g(x)$
is the right side above. Write $x=e^{-t}$ with $t>0$. For $y=x^n$ with
$n>N$ one has $-\log(1-y)\le y/(1-y)\le x^n/(1-x)$. Summing over the
primes $p>N$ gives
\begin{equation}\label{eq:PNphi}
  e^{-\Phi(t)}\ \le\ P_N(e^{-t})\ \le\
  \exp\Bigl(-\Phi(t)+\frac{e^{-(N+1)t}}{(1-e^{-t})^2}\Bigr).
\end{equation}
Put $\ell=\log N\ge1$ and $s=2\ell/N$. If $t\ge s$, then $t\ge2/N$ and
$1-e^{-t}\ge t/(1+t)\ge2/(N+2)\ge1/N$, while
$e^{-(N+1)t}\le e^{-Nt}\le N^{-2}$. So the error term in
\eqref{eq:PNphi} is at most $1$, and
\[
  g(e^{-t})\le e^{1-\Phi(t)}\cdot Ne^{-Nt}\le eNe^{-M_N}.
\]
If $0<t<s$, then $P_N(e^{-t})\le P_N(e^{-s})$ because $P_N$ is
nonincreasing, so $g(e^{-t})\le e^{1-\Phi(s)}$. Moreover
$\Phi(s)\ge M_N-Ns=M_N-2\ell$, hence $g(e^{-t})\le eN^2e^{-M_N}$. This
proves (a).

For (b), let $M\ge q$. Since $\Prm_N\cup\{q\}\subseteq\Prm_M$,
$P_N-P_M\ge P_N-P_{\Prm_N\cup\{q\}}=x^qP_N\ge0$ on $[0,1]$. On the window
$x=e^{-t}$ with $\tau\le t\le(1+\varepsilon)\tau$, \eqref{eq:PNphi}
gives $x^qP_N(x)\ge e^{-(1+\varepsilon)q\tau-B}$. The window has length
$e^{-\tau}(1-e^{-\varepsilon\tau})\ge e^{-\tau}\varepsilon\tau/(1+\varepsilon\tau)$.
Integrating and letting $M\to\infty$ proves (b).
\end{proof}

\begin{proof}[Proof of Theorem~\ref{thm:tail}]
Let $\ell=\log N$ and $\lambda_N=2\sqrt{2aN/\ell}=2\pi\sqrt{N/(3\ell)}$.
Then $\ell=o(\lambda_N)$ and $\lambda_N=o(N)$. Fix
$\varepsilon\in(0,\tfrac14)$. By Lemma~\ref{lem:phi} there is
$t_0\in(0,1)$ with
\[
  (1-\varepsilon)a\le tL(t)\Phi(t)\le(1+\varepsilon)a
  \qquad(0<t<t_0).
\]

\emph{Lower bound for $-\log(I_N-\Lambda)$.} We bound $Nt+\Phi(t)$ from
below for $t>0$. If $t\ge t_0$, then $Nt+\Phi(t)\ge Nt_0$. If $t<t_0$
and $L(t)\le(\tfrac12+\varepsilon)\ell$, then $\Phi(t)\ge b/t$ with
$b=(1-\varepsilon)a/((\tfrac12+\varepsilon)\ell)$, so
\[
  Nt+\Phi(t)\ \ge\ 2\sqrt{bN}
  =\sqrt{\frac{1-\varepsilon}{1+2\varepsilon}}\,\lambda_N .
\]
Otherwise $t<t_1=N^{-1/2-\varepsilon}$. Since $\Phi$ is decreasing,
$\Phi(t)\ge\Phi(t_1)\ge(1-\varepsilon)aN^{1/2+\varepsilon}/((\tfrac12+\varepsilon)\ell)$,
provided $t_1<t_0$. As $Nt_0$ and $N^{1/2+\varepsilon}/\ell$ are both
$\gg\lambda_N$, we get $M_N\ge\sqrt{(1-\varepsilon)/(1+2\varepsilon)}\,\lambda_N$
for large $N$. Lemma~\ref{lem:sandwich}(a) then gives
\[
  -\log(I_N-\Lambda)\ \ge\ M_N-2\ell-1
  \ \ge\ \Bigl(\sqrt{\tfrac{1-\varepsilon}{1+2\varepsilon}}-o(1)\Bigr)\lambda_N .
\]

\emph{Upper bound.} By the prime number theorem, for large $N$ there is
a prime $q$ with $N<q\le(1+\varepsilon)N$. Let $\tau=\sqrt{2a/(N\ell)}$,
so that $N\tau=\lambda_N/2$ and $\tau\to0$. For
$\tau\le t\le(1+\varepsilon)\tau$ and large $N$, we have $t<t_0$ and
\[
  L(t)\ \ge\ L(\tau)-\log(1+\varepsilon)
  =\tfrac12\bigl(\ell+\log\ell-\log2a\bigr)-\log(1+\varepsilon)
  \ \ge\ (1-\varepsilon)\tfrac{\ell}{2}.
\]
Hence
$\Phi(t)\le(1+\varepsilon)a/(\tau(1-\varepsilon)\ell/2)
=\tfrac{1+\varepsilon}{1-\varepsilon}N\tau$. By
Lemma~\ref{lem:sandwich}(b), with $(1+\varepsilon)q\tau\le(1+\varepsilon)^2N\tau$
and $\log(1/\tau)\le\ell$,
\[
  -\log(I_N-\Lambda)\ \le\
  \Bigl(\frac{1+\varepsilon}{1-\varepsilon}+(1+\varepsilon)^2\Bigr)
  \frac{\lambda_N}{2}+\ell+O_\varepsilon(1).
\]
Both constants tend to $1$ as $\varepsilon\downarrow0$, which proves the
first statement.

For the second, the prime number theorem gives $p_n\sim n\log n$, hence
$\log p_n\sim\log n$ and $p_n/\log p_n\sim n$, so
$\lambda_{p_n}\sim(2\pi/\sqrt3)\sqrt n$.
\end{proof}

\begin{remark}[Chebyshev bounds]\label{rem:chebyshev}
With only $\theta(x)\asymp x$ in place of the prime number theorem, the
same proof gives $-\log(I_N-\Lambda)\asymp\sqrt{N/\log N}$. Indeed,
Lemmas~\ref{lem:expsums} and~\ref{lem:phi} hold with $\asymp$ in place of
limits, and Bertrand's postulate supplies $q\in(N,2N]$. The prime number
theorem is what fixes the constant $2\pi/\sqrt3=2\sqrt{2\zeta(2)}$.
\end{remark}

\subsection{The next prime}\label{sec:nextprime}
Since $\Prm_N$ changes only at primes, $I_N=I_{q_N-1}$, where $q_N$ is
the least prime exceeding $N$. The two halves of
Lemma~\ref{lem:sandwich} are therefore sharpest when evaluated at
$q_N$.

\begin{proposition}[Next-prime sandwich]\label{prop:nextprime}
For every $N\ge3$,
\[
  M_{q_N}-2\log q_N-3\ \le\ -\log(I_N-\Lambda)\ \le\ M_{q_N}+\log(q_N+1)+4 .
\]
\end{proposition}

\begin{proof}
Write $q=q_N\ge5$. We first record $\Phi(1)<1$: indeed
$\Phi(1)\le\sum_{n\ge2}e^{-n}/(1-e^{-n})\le e^{-2}/((1-e^{-1})(1-e^{-2}))<1$.
Hence $M_q\le q+\Phi(1)<q+1$.

\emph{Lower bound.} Lemma~\ref{lem:sandwich}(a) with $q-1\ge3$ in place
of $N$ gives
$I_N-\Lambda=I_{q-1}-\Lambda\le e(q-1)^2e^{-M_{q-1}}$. Since $\Phi\ge0$,
$(q-1)t+\Phi(t)\ge(1-1/q)(qt+\Phi(t))$ for all $t>0$, so
$M_{q-1}\ge M_q-M_q/q>M_q-2$. Therefore
$-\log(I_N-\Lambda)\ge M_q-2-1-2\log(q-1)$.

\emph{Upper bound.} Choose $\tau>0$ with $q\tau+\Phi(\tau)\le M_q+1$;
then $q\tau<M_q+1<q+2$, so $\tau<2$. Apply Lemma~\ref{lem:sandwich}(b)
with this $q$ and $\tau$, with $\varepsilon=1/(q\tau)$, and with
$B=\Phi(\tau)$, which is admissible because $\Phi$ is decreasing. Since
$\varepsilon\tau=1/q$,
\[
  -\log(I_N-\Lambda)\le\Phi(\tau)+q\tau+1+\tau+\log(q+1)
  \le M_q+\log(q+1)+4. \qedhere
\]
\end{proof}

The proposition uses neither the prime number theorem nor any
information about the size of $q_N-N$. To compare $M_{q_N}$ with $M_N$
we need the shape of $M$. The function $\Phi$ is $C^1$, strictly
convex and strictly decreasing on $(0,\infty)$, since each summand
$-\log(1-e^{-pt})$ is, and $\Phi(t)\to\infty$ as $t\downarrow0$.
Hence, for every real $N>0$, $Nt+\Phi(t)$ has a unique minimizer $t_N$,
characterized by $-\Phi'(t_N)=N$. As an infimum of affine functions of
$N$, $M$ is concave, and it has derivative $M'(N)=t_N$, which decreases
in $N$. Consequently
\begin{equation}\label{eq:Mslope}
  (q-N)\,t_q\ \le\ M_q-M_N\ \le\ (q-N)\,t_N\qquad(N\le q).
\end{equation}
By convexity,
$(\Phi(t)-\Phi(2t))/t\le-\Phi'(t)\le(\Phi(t/2)-\Phi(t))/(t/2)$, and
Lemma~\ref{lem:phi} evaluates both sides. The result is
$-\Phi'(t)\asymp1/(t^2L(t))$ as $t\downarrow0$, with constants $a/2$
and $2a$ up to factors $1+o(1)$. Since $-\Phi'(t_N)=N$, this gives
$t_N\asymp(N\log N)^{-1/2}$.

\begin{corollary}[Second-order term and prime gaps]\label{cor:gaps}
The following are equivalent:
\begin{enumerate}
\item[(i)] $-\log(I_N-\Lambda)\le M_N+O(\log N)$ for all $N\ge3$;
\item[(ii)] $M_{q_N}-M_N=O(\log N)$;
\item[(iii)] $p_{n+1}-p_n=O\bigl(\sqrt{p_n}\,(\log p_n)^{3/2}\bigr)$.
\end{enumerate}
In particular they hold under the Riemann hypothesis. Unconditionally,
$-\log(I_N-\Lambda)=M_N+O(N^{1/40})$.
\end{corollary}

\begin{proof}
By Bertrand's postulate $q_N\le2N$, so $\log q_N=O(\log N)$, and
$M_{q_N}\ge M_N$. Then (i)$\Rightarrow$(ii) follows from the lower bound
of Proposition~\ref{prop:nextprime}, and (ii)$\Rightarrow$(i) from its
upper bound. Since $t_q\asymp t_N\asymp(N\log N)^{-1/2}$ for
$N\le q\le2N$, \eqref{eq:Mslope} gives
$M_{q_N}-M_N\asymp(q_N-N)(N\log N)^{-1/2}$. So (ii) is equivalent to
$q_N-N=O(\sqrt N(\log N)^{3/2})$. This bound for all $N$ is equivalent
to (iii): take $N=p_n$ for one direction, and use $q_N-N\le p_{n+1}-p_n$
for $p_n\le N<p_{n+1}$ for the other. Under the Riemann hypothesis
Cram\'er~\cite{Cramer1920} proved $p_{n+1}-p_n=O(\sqrt{p_n}\log p_n)$.
Unconditionally, Baker, Harman and Pintz~\cite{BHP2001} proved
$p_{n+1}-p_n=O(p_n^{0.525})$, so
$M_{q_N}-M_N=O(N^{0.525}(N\log N)^{-1/2})=O(N^{1/40})$; combine this with
Proposition~\ref{prop:nextprime}.
\end{proof}

Thus refining Theorem~\ref{thm:tail} to additive precision $O(\log N)$
is a statement about gaps between primes, and it is out of reach
unconditionally. The truncation error sees the gap after $N$ directly.
For instance, at the record gap $31397\to31469$, $I_N$ stays constant
for $72$ consecutive values of $N$ while $M_N$ increases.

\subsection{Stability of the minimizer}\label{sec:stability}
How much must $F_c$ increase when a primitive set omits a prime? The
relative theorem answers this exactly.

\begin{proposition}[Omitting a prime]\label{prop:omit}
Let $p$ be prime and $B_p=(\Prm\setminus\{p\})\cup\{p^2\}$. For every
$c\in(0,1]$, $B_p$ is the unique minimizer of $F_c$ among primitive sets
$A$ with $p\notin A$, and
\[
  F_c(B_p)=\mathcal L(c)+D_c(p),\qquad
  D_c(p)=\sum_{j=1}^{p-1}\int_0^cx^{jp}Q(x)\,dx\ \ge\
  \Delta_c(p):=\int_0^cx^pQ(x)\,dx>0,
\]
with $D_c(2)=\Delta_c(2)$.
\end{proposition}

\begin{proof}
A primitive set avoids $p$ exactly when it lies in
$U=\{2,3,\dots\}\setminus\{p\}$, and $\operatorname{Min}(U)=B_p$: a
composite with a prime factor $q\ne p$ has the proper divisor $q\in U$;
$p^k$ with $k\ge3$ has the proper divisor $p^2\in U$; primes $q\ne p$ and
$p^2$ have no proper divisor in $U$. Theorem~\ref{thm:ambient} gives the
first statement. On $[0,1)$, $Q=(1-x^p)P_{\Prm\setminus\{p\}}$ and
$P_{B_p}=(1-x^{p^2})P_{\Prm\setminus\{p\}}$, so
$P_{B_p}=Q\cdot(1-x^{p^2})/(1-x^p)=Q\sum_{j=0}^{p-1}x^{jp}$; integrate.
\end{proof}

\begin{corollary}[Stability]\label{cor:stability}
Let $A$ and $A_1,A_2,\dots$ be primitive sets.
\begin{enumerate}
\item[(a)] If $p\notin A$ then $F_c(A)-\mathcal L(c)\ge D_c(p)\ge\Delta_c(p)$
  for every $c\in(0,1]$.
\item[(b)] If $F_1(A)-\Lambda<\Delta_1(p)$ for every prime $p\le R$, then
  $A\cap[2,R]=\Prm_R$. Since $\Delta_1$ is decreasing in $p$, it suffices
  that $F_1(A)-\Lambda<\Delta_1(r)$ for the largest prime $r\le R$.
\item[(c)] $F_1(A_j)\to\Lambda$ if and only if every prime belongs to $A_j$
  for all large $j$.
\end{enumerate}
\end{corollary}

\begin{proof}
(a) is Proposition~\ref{prop:omit}. For (b), (a) shows that every prime
$p\le R$ lies in $A$, and then no composite $m\le R$ lies in $A$, since
its prime factors are at most $R$. For (c), if some prime is omitted
infinitely often then (a) bounds $F_1(A_j)-\Lambda$ below along a
subsequence. Conversely, if every prime is eventually in $A_j$, then by
primitivity every fixed composite is eventually excluded; the indicator
of $A_j$ converges pointwise to that of $\Prm$, so $P_{A_j}\to Q$ on
$[0,1)$ by dominated convergence for the products, and $F_1(A_j)\to\Lambda$
by dominated convergence for the integrals.
\end{proof}

The threshold $\Delta_1(p)$ is governed by the same saddle as the
truncation error.

\begin{proposition}[Size of the threshold]\label{prop:delta}
For every prime, and indeed every integer, $p\ge2$,
\[
  M_p\ \le\ -\log\Delta_1(p)\ \le\ M_p+\log p+3,
\]
and consequently $-\log\Delta_1(p)\sim2\pi\sqrt{p/(3\log p)}$.
\end{proposition}

\begin{proof}
With $x=e^{-t}$, $x^pQ(x)=e^{-(pt+\Phi(t))}\le e^{-M_p}$, which gives the
lower bound after integrating over $(0,1)$. For the upper bound, fix
$\tau>0$ and integrate over the window $t\in[\tau,\tau+1/p]$, where
$pt+\Phi(t)\le p\tau+\Phi(\tau)+1$ because $\Phi$ is decreasing. The
window has $x$-length $e^{-\tau}(1-e^{-1/p})\ge e^{-\tau}/(p+1)$, so
\[
  -\log\Delta_1(p)\le p\tau+\Phi(\tau)+\tau+1+\log(p+1)
  \le p\tau+\Phi(\tau)+\tau+\frac1p+1+\log p .
\]
Choose $\tau$ with $p\tau+\Phi(\tau)\le M_p+\delta$, where
$\delta=(1-\Phi(1))/3>0$ (recall $\Phi(1)<1$ from the proof of
Proposition~\ref{prop:nextprime}). Since $M_p\le p+\Phi(1)$, this gives
$\tau\le1+(\Phi(1)+\delta)/p$, and the bound becomes
$M_p+\log p+2+\delta+(\Phi(1)+\delta+1)/p\le M_p+\log p+3$ for $p\ge2$.
For the asymptotic, the proof of Theorem~\ref{thm:tail} gives
$M_p\ge(1-o(1))\lambda_p$. Lemma~\ref{lem:sandwich}(a) gives
$M_p\le-\log(I_p-\Lambda)+2\log p+1$, and Theorem~\ref{thm:tail} gives
$-\log(I_p-\Lambda)\sim\lambda_p$.
\end{proof}

So a primitive set whose $F_1$ lies within
$e^{-(1+o(1))2\pi\sqrt{R/(3\log R)}}$ of $\Lambda$ agrees with the primes
up to $R$ (Corollary~\ref{cor:stability}(b) with the largest prime
$r\le R$, and $r\sim R$). This scale is sharp: the truncation
$\Prm_{r-1}$ omits $r$, and its excess is
$I_{r-1}-\Lambda=e^{-(1+o(1))\lambda_r}$ by Theorem~\ref{thm:tail}.

\subsection{Second-order terms}
Relative refinements of Theorem~\ref{thm:tail} do not need short gaps.
Put $a=\zeta(2)$, $\ell=\log N$, $\lambda_N=2\sqrt{2aN/\ell}
=2\pi\sqrt{N/(3\ell)}$,
\[
  \sigma=\sum_{m\ge1}\frac{\log m}{m^2}=-\zeta'(2),\qquad
  \kappa=\gamma+\frac{\sigma}{a},\qquad
  C=\kappa+\tfrac12\log(2a)=1.7426003995\ldots,
\]
where $\gamma$ is Euler's constant.

\begin{theorem}[Second-order tail asymptotics]\label{thm:second}
As $N\to\infty$,
\[
  -\log(I_N-\Lambda)=\lambda_N\Bigl(1-\frac{\log\ell}{2\ell}+\frac{C}{\ell}
  +O\Bigl(\frac{(\log\ell)^2}{\ell^2}\Bigr)\Bigr).
\]
\end{theorem}

The only arithmetic input is the prime number theorem with the error
term $\theta(x)=x+O(x/\log^2x)$, which follows from the classical
de la Vall\'ee Poussin error term. The proof has four steps.

\emph{Step 1: prime sums.} Put $L=\log(1/u)$. Then
\[
  S(u)=\frac1{uL}\Bigl(1+\frac{\gamma}{L}\Bigr)+O\Bigl(\frac1{uL^3}\Bigr)
  \qquad(u\downarrow0).
\]
Partial summation gives
$S(u)=\sum_{n\ge2}\theta(n)\bigl(g(n)-g(n+1)\bigr)$ with
$g(x)=e^{-xu}/\log x$. The error term of $\theta$ contributes
$O(\sum_n(e^{-nu}/\log^4n+nu\,e^{-nu}/\log^3n))=O(1/(uL^3))$. The main
term equals $\sum_{n\ge2}g(n)+O(1)=\int_2^\infty g+O(1)$. Finally,
$u\int_2^\infty g=\int_{2u}^\infty e^{-y}\,dy/(L+\log y)$; expanding
$1/(L+\log y)$ to second order and using
$\int_0^\infty e^{-y}\log y\,dy=\Gamma'(1)=-\gamma$ gives the claim.

\emph{Step 2: the exponent.} Put $L=\log(1/t)$. Then
\[
  \Phi(t)=\frac{a}{tL}\Bigl(1+\frac{\kappa}{L}\Bigr)+O\Bigl(\frac1{tL^3}\Bigr)
  \qquad(t\downarrow0).
\]
Indeed $\Phi(t)=\sum_mS(mt)/m$ (proof of Lemma~\ref{lem:phi}). The terms
with $m>t^{-1/2}$ contribute $O(t^{-1/2})$, because $S(u)\le1/u$. For
$m\le t^{-1/2}$, Step~1 applies with $L-\log m\ge L/2$ in place of $L$,
and expanding $1/(L-\log m)$ in powers of $1/L$ gives
$\sum_m m^{-2}\bigl(1/L+(\gamma+\log m)/L^2\bigr)/t$ up to
$O(1/(tL^3))$. Here $\sum m^{-2}=a$ and $\sum m^{-2}\log m=\sigma$.

\emph{Step 3: the saddle.} Let $u_0=\frac12(\ell+\log\ell-\log2a)$, so
$e^{2u_0}=N\ell/(2a)$, and $\tau=e^{-u_0}$; then $N\tau=\lambda_N/2$. Put
\[
  S_N=1-\frac{\log\ell}{2\ell}+\frac{C}{\ell},\qquad
  r_N=\frac{\log\ell}{\ell}.
\]
Since $\Phi$ is nonnegative and nonincreasing, Step~2 and
Lemma~\ref{lem:saddle} (Appendix~\ref{app:saddle}) give a constant $K$
such that, for all large $N$,
\[
  \lambda_N\bigl(S_N-Kr_N^2\bigr)\le Nt+\Phi(t)\quad(t>0),\qquad
  N\tau+\Phi(\tau)\le\lambda_N\bigl(S_N+Kr_N^2\bigr).
\]
In particular $M_N=\inf_{t>0}(Nt+\Phi(t))=\lambda_N\bigl(S_N+O(r_N^2)\bigr)$.

\emph{Step 4: transfer.} Lemma~\ref{lem:sandwich}(a) gives
$-\log(I_N-\Lambda)\ge M_N-2\log N-1$. For the upper bound, the error
term of $\theta$ provides a prime $q\in(N,N+O(N/\ell^2)]$. Apply
Lemma~\ref{lem:sandwich}(b) with $\tau=e^{-u_0}$, $\varepsilon=\ell^{-2}$
and $B=\Phi(\tau)$, which is admissible because $\Phi$ is decreasing:
\[
  -\log(I_N-\Lambda)\le N\tau+\Phi(\tau)+(q(1+\varepsilon)-N)\tau
  +O(\log N),
\]
and $(q(1+\varepsilon)-N)\tau=O(N\tau/\ell^2)=O(\lambda_N/\ell^2)$.
Since $\log N=o(\lambda_N(\log\ell/\ell)^2)$, the theorem follows.

\begin{remark}[Partitions into primes]\label{rem:partitions}
Let $p_{\Prm}(N)$ be the number of partitions of $N$ into primes, so
that $\sum_Np_{\Prm}(N)x^N=1/Q(x)=e^{\Phi(t)}$ with $x=e^{-t}$. By Roth
and Szekeres~\cite{RothSzekeres},
$\log p_{\Prm}(N)\sim2\pi\sqrt{N/(3\log N)}$, so Theorem~\ref{thm:tail}
says $-\log(I_N-\Lambda)\sim\log p_{\Prm}(N)$: both are governed by the
saddle point of $Nt+\Phi(t)$, although the proofs above do not use
partitions. More precisely,
\begin{equation}\label{eq:pP}
  \log p_{\Prm}(N)=M_N+O(\log N).
\end{equation}
The upper bound $\log p_{\Prm}(N)\le M_N$ holds for every $N$, since
$p_{\Prm}(N)e^{-Nt}\le e^{\Phi(t)}$. For the lower bound, Vaughan's
asymptotic formula~\cite{Vaughan2008}, in the saddle-point form of
Gafni~\cite[Theorem~1.3 with $k=1$]{Gafni2021}, gives
$p_{\Prm}(N)=e^{M_N}\bigl(2\pi\Phi''(t_N)\bigr)^{-1/2}\bigl(1+O(N^{-1/5})\bigr)$,
with $t_N$ as in Section~\ref{sec:nextprime}, and
$\log\Phi''(t_N)=O(\log N)$: the term $p=2$ gives
$-\Phi'(t)\ge1/t-1$, so $t_N\ge1/(N+1)$, while $Nt_N\le M_N\le N+\Phi(1)<N+1$ (proof of Proposition~\ref{prop:nextprime}) gives
$t_N\le2$; on $(0,2]$,
$4e^{-4}\le\Phi''(t)\le2(1+t)^5t^{-5}$, using
$1-e^{-pt}\ge t/(1+t)$ and $\sum_nn^2y^n\le2/(1-y)^3$. (The prefactor
contributes $-\tfrac34\log N+O(\log\log N)$~\cite[(1.5)]{Gafni2021}.)
By \eqref{eq:pP} and Corollary~\ref{cor:gaps}, the difference
$-\log(I_N-\Lambda)-\log p_{\Prm}(N)$ is governed by the gap after $N$.
Since $\log N=o(\lambda_Nr_N^2)$, Step~3 above also gives
$\log p_{\Prm}(N)=\lambda_N\bigl(S_N+O(r_N^2)\bigr)$. This agrees with
the formal saddle-point expansion of Bartel, Bhaduri, Brack and
Murthy~\cite[(43)]{BBBM2017}: their coefficient of $\log\ell/\ell$ is
$-\tfrac12$ (which they give in place of the coefficient $+1$
in~\cite{Vaughan2008}), and their coefficient of $1/\ell$ is
$f_2/f_1+\log(\pi/\sqrt3)$ with $f_1=a$ and $f_2=\gamma a+\sigma$
\cite[(23)]{BBBM2017}, which equals $\kappa+\tfrac12\log(2a)=C$.
\end{remark}

\subsection{Open problems}
\begin{enumerate}
\item Is there a closed form for $\Lambda$, as for the all-exponent and
even members of the family? Is $\Lambda$ irrational?
\item Characterize the increasing functions $\phi$ for which $\Prm_N$
minimizes $\E\phi(Z_A)$ over primitive $A\subseteq\{2,\dots,N\}$ for
every $N$ (notation of Remark~\ref{rem:icv}). This convex cone contains
the increasing concave functions and the powers $x^k$, $0<k\le17/4$,
and excludes the indicators $\mathbf 1_{(x,1]}$ for $x$ near $1$
(Proposition~\ref{prop:pointwise}). It omits $x^k$ for large $k$: for
finite $A$,
$\E Z_A^k\sim\bigl(\prod_{a\in A}a\bigr)|A|!\,k^{-|A|}$ as
$k\to\infty$, and already $\E Z_{A_9}^k<\E Z_{\Prm_9}^k$ for every
$k\ge126$ (exact computation, \texttt{powers\_check.py}). The bound
$2\alpha_0\approx4.286$ of Remark~\ref{rem:icv} is not sharp: for
$N\in\{12,18,24\}$ and every maximal primitive subset of
$\{2,\dots,N\}$, the primes give the smaller $\E\min(Z^k,c^k)$ for
eleven exponents $k\in[5/2,20]$ and
$c\in\{\tfrac12,\tfrac9{10},\tfrac{19}{20},\tfrac{99}{100},1\}$
(\texttt{moments\_search.py}; 40-digit arithmetic, not a proof). What
is the largest admissible $k$?
\item Proposition~\ref{prop:nextprime} determines $-\log(I_N-\Lambda)$
up to $O(\log N)$ by $M_{q_N}$, and Corollary~\ref{cor:gaps} shows that
replacing $M_{q_N}$ by $M_N$ at this precision is a prime-gap problem.
What is the next term relative to $M_{q_N}$? Writing $I_N-\Lambda$ as
the sum over primes $p>N$ of $\int_0^1x^pP_{\Prm_{<p}}(x)\,dx$, applying
Laplace's method to each term, and giving the primes near $q_N$ their
average density on the scale $\sqrt{N\log N}$ suggests
\[
  -\log(I_N-\Lambda)=M_{q_N}+\tfrac14\log q_N+\tfrac34\log\log q_N
  -\tfrac12\log\pi+\tfrac14\log(2\zeta(2))+o(1).
\]
Numerically, at $N=p_n$ for $20\le n\le9000$ ($N\le93{,}000$, including
the maximal gaps after $19609$ and $31397$), the left side minus the
right side, without the $o(1)$, lies between $-0.23$ and $0.11$ (script
\texttt{next\_prime\_check.py}; floating point, not certified). A proof
for every $N$ would need control of primes in intervals of length about
$\sqrt N$ after $N$, beyond current unconditional methods. Does the
formula hold for almost all $N$?
\end{enumerate}

\section{Formal verification}\label{sec:formal}

The accompanying Lean~4/Mathlib development~\cite{Lean4,Mathlib}
(\url{https://github.com/alerad/prime-extremality}, release \texttt{v1.0.1}, commit \texttt{f6caaa43},~\cite{Artifact}) is built with Lean
\texttt{v4.35.0}, Mathlib commit
\texttt{56a00585\allowbreak 5ea17e0b\allowbreak 0958f73e\allowbreak c801ae14\allowbreak 0b0a1679},
the \href{https://github.com/alerad/leancert}{LeanCert library}
(Bernstein certificates, interval arithmetic and exact $q$-product
integrals), and the prime number theorem imported,
unmodified, from the PrimeNumberTheoremAnd project~\cite{PNTplus} at
commit \texttt{127b50e5}. It formalizes, for every $c\in(0,1]$ and for
finite and infinite sets: Theorems~\ref{thm:unique},
\ref{thm:profile}, \ref{thm:running} (by the route of
Section~\ref{sec:running}) and~\ref{thm:ambient};
Propositions~\ref{prop:pointwise} and~\ref{prop:block}; the
increasing-concave minimality of Remark~\ref{rem:icv} for finite sets
and its moment extension for all real $0<k\le17/4$, including the exact
values of $J_{17/8}(1)$ and $J_3(1)$; the $26$ digits of $\Lambda$;
Lemmas~\ref{lem:expsums}, \ref{lem:phi}, \ref{lem:sandwich}
and~\ref{lem:saddle}; Theorems~\ref{thm:tail} and~\ref{thm:second};
Proposition~\ref{prop:omit}, Corollary~\ref{cor:stability} and
Proposition~\ref{prop:delta}; and the comparison at $k=126$ in open
problem~2. The prime number theorem enters through
\texttt{chebyshev\_asymptotic} and \texttt{nth\_prime\_asymp} for
Theorem~\ref{thm:tail}, and through \texttt{MediumPNT} together with
Mathlib's bound for $\psi-\theta$ for the hypothesis
$\theta(x)=x+O(x/\log^2x)$ of Theorem~\ref{thm:second}. Not formalized
are the probabilistic model of Remark~\ref{rem:icv} itself, the infinite
case of its increasing-concave statement, the numerical value of
$\alpha_0$, the identification of $\operatorname{Min}(U)$ in
Corollary~\ref{cor:ambient}, Proposition~\ref{prop:nextprime},
Corollary~\ref{cor:gaps}, and Remarks~\ref{rem:chebyshev},
\ref{rem:partitions} and~\ref{rem:additive}. Some Lean proofs take
equivalent routes (for instance dominated convergence in place of
Stieltjes integration by parts); the repository lists each deviation and
maps every result to its Lean declaration. An earlier
artifact~\cite{ProductInvariants} verified a Bernstein-certificate proof
of Theorem~\ref{thm:running}.

All of these results depend only on the axioms \texttt{propext},
\texttt{Classical.choice} and \texttt{Quot.sound}, except the digits of
$\Lambda$ and the comparison in open problem~2, whose finite rational
computations use \texttt{native\_decide} and hence trust Lean's
compiler.

\appendix
\section{The saddle estimate}\label{app:saddle}

This appendix proves the estimate used in Step~3 of the proof of
Theorem~\ref{thm:second}. The lemma is stated for an abstract function
$\Phi$; its Lean counterpart is \texttt{saddle\_second\_order} in
\texttt{Tail/Saddle2.lean}, and the proof below follows it.

\begin{lemma}[Second-order saddle estimate]\label{lem:saddle}
Let $a>0$, $\kappa\in\mathbb R$, $K_1\ge0$ and $t_0\in(0,1)$, and let
$\Phi:(0,\infty)\to[0,\infty)$ be nonincreasing with
\begin{equation}\label{eq:saddle-hyp}
  \Bigl|\Phi(t)-\frac{a}{tL}\Bigl(1+\frac{\kappa}{L}\Bigr)\Bigr|
  \le\frac{K_1}{tL^3}\qquad(0<t<t_0,\ L=\log(1/t)).
\end{equation}
For $N>e$ put $\ell=\log N$, $\lambda_N=2\sqrt{2aN/\ell}$,
$u_0=\frac12(\ell+\log\ell-\log2a)$, $\tau=e^{-u_0}$,
$C=\kappa+\frac12\log(2a)$,
$S_N=1-\log\ell/(2\ell)+C/\ell$ and $r_N=\log\ell/\ell$. There is a
constant $K$, depending only on $a,\kappa,K_1$, such that for all
sufficiently large $N$:
\begin{enumerate}
\item[(i)] $Nt+\Phi(t)\ge\lambda_N(S_N-Kr_N^2)$ for every $t>0$;
\item[(ii)] $N\tau+\Phi(\tau)\le\lambda_N(S_N+Kr_N^2)$, and
  $N\tau=\lambda_N/2$.
\end{enumerate}
\end{lemma}

\begin{proof}
\emph{Notation.} Write $K_0=K_1/a$ and
$h_\pm(u)=1/u+\kappa/u^2\pm K_0/u^3$. In the variable $u=\log(1/t)$,
\eqref{eq:saddle-hyp} says
\begin{equation}\label{eq:saddle-hpm}
  a e^uh_-(u)\le\Phi(e^{-u})\le a e^uh_+(u)\qquad(u>\log(1/t_0)).
\end{equation}
Put $X=e^{u_0}$ and $A=N/X$. Then $X^2=N\ell/(2a)$, so $A=2aX/\ell$,
$A^2=2aN/\ell$ and $\lambda_N=2A$. For $t=e^{-(u_0+v)}$,
\begin{equation}\label{eq:saddle-scale}
  Nt=Ae^{-v},\qquad ae^{u_0+v}=A\,\tfrac{\ell}{2}\,e^{v}.
\end{equation}
Let $r=r_N$, $D=1+|\log2a|$, $x=(2u_0-\ell)/\ell=(\log\ell-\log2a)/\ell$
and $p=\ell/(2u_0)=1/(1+x)$. Then $S_N=1-x/2+\kappa/\ell$. We take $N$
so large that $\log\ell\ge1$ and
\begin{equation}\label{eq:saddle-large}
  (D+|\kappa|+2)\,r\le\tfrac1{40},\qquad
  u_0\ge U:=20(|\kappa|+K_0)+1,\qquad Xt_1\ge4,
\end{equation}
where $t_1=\min(t_0,e^{-U})$. Then $1/\ell\le r\le1/80$, so $\ell\ge80$;
$|x|\le Dr\le1/40$, so $p\le2$; and $S_N\le\frac65$.

\emph{Algebra at $u_0$.} Put $P_\pm=\frac\ell2h_\pm(u_0)
=p+2\kappa p^2/\ell\pm4K_0p^3/\ell^2$. Using $p=1-x+x^2p$ and
$p^2-1=-xp(p+1)$,
\[
  1+P_\pm-2S_N=x^2p+\frac{2\kappa(p^2-1)}{\ell}\pm\frac{4K_0p^3}{\ell^2},
\]
so that
\begin{equation}\label{eq:saddle-alg}
  |1+P_\pm-2S_N|\le K_2r^2,\qquad K_2=2D^2+12|\kappa|D+32K_0 .
\end{equation}
Similarly $\delta=P_--1=(p-1)+2\kappa p^2/\ell-4K_0p^3/\ell^2$
satisfies $|\delta|\le K_3r$ with $K_3=2D+8|\kappa|+32K_0$, because
$|p-1|=|x|p$.

\emph{Two elementary inequalities.} First, $e^{-v}+e^v\ge2+v^2$ for
all real $v$, and $|e^v-1|\le2|v|$ for $v\le1$ (for $0\le v\le1$,
$e^v-1\le v(e-1)$; for $v<0$, $1-e^v\le\min(1,|v|)$). Hence for $v\le1$,
$\delta\in\mathbb R$ and $B\ge0$,
\begin{equation}\label{eq:saddle-quad}
  e^{-v}+(1+\delta)e^v-B\max(v,0)
  \ge2+\delta+\tfrac{v^2}2-(2|\delta|+B)|v|
  \ge2+\delta-\tfrac12(2|\delta|+B)^2 .
\end{equation}
Second, $h_-$ is nonincreasing on $[U,\infty)$, since
$h_-'(u)=-u^{-2}(1+2\kappa/u-3K_0/u^2)$ and
$(2|\kappa|+3K_0)/u<1$ there. Moreover, for $u_0>0$ and $0\le v\le1$,
\begin{equation}\label{eq:saddle-step}
  h_-(u_0+v)\ge h_-(u_0)-\frac{v}{u_0^2}-\frac{2|\kappa|}{u_0^3},
\end{equation}
because $1/(u_0+v)\ge1/u_0-v/u_0^2$,
$0\le u_0^{-2}-(u_0+v)^{-2}\le2u_0^{-3}$ and
$(u_0+v)^{-3}\le u_0^{-3}$.

\emph{Proof of (i).} Fix $t>0$.

\emph{Case $t\ge t_1$.} Since $\Phi\ge0$,
$Nt+\Phi(t)\ge Nt_1=AXt_1\ge4A=2\lambda_N\ge\lambda_NS_N$.

\emph{Case $t<t_1$.} Write $t=e^{-(u_0+v)}$; then $u_0+v>U$ and
$t<t_0$, so \eqref{eq:saddle-hpm} and \eqref{eq:saddle-scale} give
\[
  Nt+\Phi(t)\ge A\Bigl(e^{-v}+\tfrac\ell2e^vh_-(u_0+v)\Bigr).
\]
If $v\le0$, then $U<u_0+v\le u_0$ and monotonicity of $h_-$ gives
$\frac\ell2h_-(u_0+v)\ge P_-$, so by \eqref{eq:saddle-quad} with $B=0$,
$Nt+\Phi(t)\ge A(2+\delta-2\delta^2)$. If $0<v\le1$, then
\eqref{eq:saddle-step} and $e^v\le3$ give
$\frac\ell2e^vh_-(u_0+v)\ge P_-e^v-Bv-E$ with
$B=3\ell/(2u_0^2)=6p^2/\ell\le24r$ and
$E=3|\kappa|\ell/u_0^3=24|\kappa|p^3/\ell^2\le192|\kappa|r^2$; by
\eqref{eq:saddle-quad},
$Nt+\Phi(t)\ge A\bigl(2+\delta-\frac12(2|\delta|+B)^2-E\bigr)$. In both
subcases, by \eqref{eq:saddle-alg} and $2+\delta=1+P_-$,
\[
  Nt+\Phi(t)\ge A\Bigl(2S_N-K_2r^2-\tfrac12(2K_3+24)^2r^2
  -192|\kappa|r^2\Bigr)=\lambda_N(S_N-Kr^2),
\]
with $K=\frac12K_2+\frac14(2K_3+24)^2+96|\kappa|$. Finally, if $v>1$,
put $t_2=e^{-(u_0+1)}<t_1$. Since $\Phi$ is nonincreasing,
\eqref{eq:saddle-hpm} and \eqref{eq:saddle-scale} give
\[
  \Phi(t)\ge\Phi(t_2)\ge Ae\,\tfrac\ell2h_-(u_0+1)
  =Ae\,\frac{\ell}{2(u_0+1)}\bigl(1+\kappa q-K_0q^2\bigr),
  \qquad q=\frac1{u_0+1}.
\]
Here $\ell/(2(u_0+1))=\ell/(\ell(1+x)+2)\ge\frac{19}{20}$ because
$x\le\frac1{40}$ and $\ell\ge80$, and
$1+\kappa q-K_0q^2\ge1-(|\kappa|+K_0)q\ge\frac{19}{20}$ because
$u_0+1\ge U$. Hence $\Phi(t)\ge(\frac{19}{20})^2eA>2.4A\ge\lambda_NS_N$.
This proves (i), since $S_N\ge S_N-Kr^2$.

\emph{Proof of (ii).} By \eqref{eq:saddle-large},
$\tau=1/X\le t_1/4<t_0$, and $N\tau=A=\lambda_N/2$. By
\eqref{eq:saddle-hpm} and \eqref{eq:saddle-scale} with $v=0$,
$\Phi(\tau)\le aXh_+(u_0)=AP_+$, so by \eqref{eq:saddle-alg}
\[
  N\tau+\Phi(\tau)\le A(1+P_+)\le A(2S_N+K_2r^2)
  \le\lambda_N(S_N+Kr^2). \qedhere
\]
\end{proof}

\section*{Declarations}
\emph{Use of AI.} The author used Claude Code (Anthropic) in
developing parts of the Lean code, the verification scripts, and the
manuscript. The author checked the mathematical arguments and assumes
responsibility for the paper.

\emph{Competing interests.} The author declares no competing
interests.

{\sloppy\emph{Data availability.} The Lean development of
Section~\ref{sec:formal} and the supplementary scripts
\texttt{crossing\_endpoint.py}, \texttt{powers\_check.py},
\texttt{moment\_certificate.py}, \texttt{moments\_search.py},
\texttt{tail\_\allowbreak asymptotic\_\allowbreak check.py} and
\texttt{next\_\allowbreak prime\_\allowbreak check.py} are publicly
available at \url{https://github.com/alerad/prime-extremality} and archived at \url{https://doi.org/10.5281/zenodo.23537236}~\cite{Artifact}.\par}

\begin{thebibliography}{99}

\bibitem{Alexeev2026}
B.~Alexeev, K.~Barreto, Y.~Li, J.~D. Lichtman, L.~Price, J.~I. Shah,
Q.~Tang and T.~Tao,
Primitive sets and von Mangoldt chains: Erd\H{o}s Problem \#1196 and
beyond, arXiv:2605.00301 (2026),
\url{https://arxiv.org/abs/2605.00301}.

\bibitem{BHP2001}
R.~C. Baker, G.~Harman and J.~Pintz,
The difference between consecutive primes, II,
\emph{Proc. London Math. Soc. (3)} \textbf{83} (2001), 532--562,
\url{https://doi.org/10.1112/plms/83.3.532}.

\bibitem{BanksMartin2013}
W.~D. Banks and G.~Martin,
Optimal primitive sets with restricted primes,
\emph{Integers} \textbf{13} (2013), Paper A69,
\url{https://math.colgate.edu/~integers/n69/n69.pdf}.

\bibitem{BBBM2017}
J.~Bartel, R.~K. Bhaduri, M.~Brack and M.~V.~N. Murthy,
On the asymptotic prime partitions of integers,
\emph{Phys. Rev. E} \textbf{95} (2017), 052108,
\url{https://doi.org/10.1103/PhysRevE.95.052108}.

\bibitem{Cramer1920}
H.~Cram\'er,
Some theorems concerning prime numbers,
\emph{Ark. Mat. Astr. Fys.} \textbf{15} (1920), no.~5, 1--33.

\bibitem{Gafni2021}
A.~Gafni,
Partitions into prime powers,
\emph{Mathematika} \textbf{67} (2021), no.~2, 468--488,
\url{https://doi.org/10.1112/mtk.12082}.

\bibitem{Lean4}
L.~de~Moura and S.~Ullrich,
The Lean 4 theorem prover and programming language,
in \emph{Automated Deduction -- CADE 28}, Lecture Notes in Comput.
Sci.\ 12699, Springer, 2021, 625--635,
\url{https://doi.org/10.1007/978-3-030-79876-5_37}.

\bibitem{Lichtman2023}
J.~D. Lichtman,
A proof of the Erd\H{o}s primitive set conjecture,
\emph{Forum Math. Pi} \textbf{11} (2023), e18,
\url{https://doi.org/10.1017/fmp.2023.16}.

\bibitem{LichtmanPomerance2019}
J.~D. Lichtman and C.~Pomerance,
The Erd\H{o}s conjecture for primitive sets,
\emph{Proc. Amer. Math. Soc. Ser. B} \textbf{6} (2019), 1--14,
\url{https://doi.org/10.1090/bproc/40}.

\bibitem{Mathlib}
The mathlib Community,
The Lean mathematical library,
in \emph{CPP 2020}, ACM, 2020, 367--381,
\url{https://doi.org/10.1145/3372885.3373824}.

\bibitem{Mulcrone1956}
T.~F. Mulcrone,
Definite integral of an infinite product (Solution to Problem 4614),
\emph{Amer. Math. Monthly} \textbf{63} (1956), no.~2, 127--128,
\url{https://doi.org/10.2307/2306451}.

\bibitem{PNTplus}
A.~Kontorovich and T.~Tao,
\emph{Prime Number Theorem and More}, Lean project,
\url{https://github.com/AlexKontorovich/PrimeNumberTheoremAnd}.

\bibitem{Artifact}
A.~Radisic,
\texttt{prime-extremality}, Lean~4 artifact,
\url{https://github.com/alerad/prime-extremality}, release
\texttt{v1.0.1}, commit \texttt{f6caaa43}; archived at
\url{https://doi.org/10.5281/zenodo.23537236}.

\bibitem{ProductInvariants}
A.~Radisic,
\texttt{productinvariants}, Lean~4 artifact,
\url{https://github.com/alerad/productinvariants}, commit
\texttt{d3b18c06}; archived at
\url{https://doi.org/10.5281/zenodo.21191624}.

\bibitem{RothSzekeres}
K.~F. Roth and G.~Szekeres,
Some asymptotic formulae in the theory of partitions,
\emph{Quart. J. Math. Oxford Ser.} \textbf{5} (1954), 241--259,
\url{https://doi.org/10.1093/qmath/5.1.241}.

\bibitem{Sandham1954}
H.~F. Sandham,
Problem 4614,
\emph{Amer. Math. Monthly} \textbf{61} (1954), no.~10, 718,
\url{https://doi.org/10.2307/2307339}.

\bibitem{ShakedShanthikumar}
M.~Shaked and J.~G. Shanthikumar,
\emph{Stochastic Orders},
Springer Series in Statistics, Springer, New York, 2007,
\url{https://doi.org/10.1007/978-0-387-34675-5}.

\bibitem{Vaughan2008}
R.~C. Vaughan,
On the number of partitions into primes,
\emph{Ramanujan J.} \textbf{15} (2008), no.~1, 109--121,
\url{https://doi.org/10.1007/s11139-007-9037-5}.

\end{thebibliography}

\medskip
\noindent 2020 \emph{Mathematics Subject Classification.} Primary
11B75, 11N05; Secondary 11P82, 06A07, 26D15, 60E15, 11Y60, 68V20.

\smallskip
\noindent\textit{Email address}: \texttt{aleloid@proton.me}

\end{document}
